%=====================================================================
% FINAL Q1 JOURNAL-READY PAPER: Electricity Demand QAIDS + TOU Vietnam
% Author: Truong Nhat Duy (UEH & Fulbright Vietnam)
% Date: November 07, 2025
%=====================================================================

\documentclass[12pt,a4paper]{article}
\usepackage[utf8]{inputenc}
\usepackage[T1]{fontenc}
\usepackage{amsmath, amssymb, amsthm}
\usepackage{booktabs}
\usepackage{graphicx}
\usepackage{natbib}
\usepackage{geometry}
\usepackage{float}
\usepackage{caption}
\usepackage{subcaption}
\usepackage{siunitx}
\usepackage{hyperref}
\usepackage{xcolor}
\usepackage{titlesec}
\usepackage{tocloft}
\usepackage{fancyhdr}
\usepackage{setspace}
\usepackage{pdflscape}
\usepackage{longtable}
\usepackage{rotating}
\usepackage{ifthen} % <== để điều kiện kiểm tra file ảnh

\onehalfspacing
\geometry{margin=1in}
\bibpunct{(}{)}{;}{a}{}{,}

% Custom formatting
\titleformat{\section}{\large\bfseries}{\thesection}{1em}{}
\titleformat{\subsection}{\bfseries}{\thesubsection}{1em}{}
\renewcommand{\cftsecleader}{\cftdotfill{\cftdotsep}}
\renewcommand{\cftsecfont}{\bfseries}

% Header
\pagestyle{fancy}
\fancyhf{}
\rhead{\thepage}
\lhead{\textit{Electricity Demand with QAIDS and TOU}}

\title{\textbf{Estimating Residential Electricity Demand in Vietnam Using QAIDS with Time-of-Use Pricing: Evidence from VHLSS 2022 and Policy Simulations for 2025--2030}}

\author{
	\textbf{Truong Nhat Duy}\thanks{Email: duy.truongnhat@ueh.edu.vn} \\[0.3em]
	\small MSc Candidate, Development Economics, University of Economics Ho Chi Minh City (UEH) \\[-0.2em]
	\small Master of Phaublic Policy, Fulbright School of Public Policy and Management \\[-0.2em]
	\small Ho Chi Minh City, Vietnam
}
\date{November 07, 2025}

\begin{document}
	
	\maketitle
	
	%=====================================================================
	% TABLE OF CONTENTS
	%=====================================================================
	\tableofcontents
	\newpage
	
	%=====================================================================
	% ABSTRACT
	%=====================================================================
	\begin{abstract}
		\noindent This study estimates Vietnam's residential electricity demand using the \textbf{Quadratic Almost Ideal Demand System (QAIDS)} augmented with \textbf{Time-of-Use (TOU) pricing} and temperature controls, based on the \textbf{Vietnam Household Living Standards Survey (VHLSS) 2022} ($N=9{,}383$). Results show a \textbf{price elasticity of $-0.62$}, \textbf{income elasticity of $0.38$} (rising to $0.50$ among the richest quintile), and a \textbf{TOU-urban interaction of $+0.18$}. A 10\% peak price increase causes a \textbf{Compensating Variation (CV) loss of VND 58,000/household/month}, with the poorest quintile bearing 2.1 times the relative burden. Dynamic simulations for 2025--2030 project \textbf{12--25\% peak reduction} and \textbf{10 million tons CO$_2$ savings} under full TOU + solar integration. The findings provide the first nationally representative evidence for Vietnam's TOU rollout and net-zero transition.
	\end{abstract}
	
	%=====================================================================
	% 1. INTRODUCTION
	%=====================================================================
	\section{Introduction}
	Vietnam's electricity demand grew 10\% annually from 2015--2022 \citep{gso2024}, with residential consumption accounting for 52\% of total load. The government introduced a 4.8\% tariff increase in 2025 (Decision 599/QD-EVN) and TOU pilots in Hanoi and Ho Chi Minh City \citep{evn2025}. However, tiered pricing induces endogeneity, and no study has jointly estimated demand systems with TOU and climate controls.
	
	This paper addresses three questions:
	\begin{enumerate}
		\item How do price, income, and TOU affect electricity expenditure shares?
		\item How do elasticities vary by income, region, and temperature?
		\item What are the welfare and environmental impacts of nationwide TOU?
	\end{enumerate}
	
	\textbf{Contributions}: (1) First QAIDS with TOU and temperature on VHLSS 2022; (2) CV/EV welfare analysis; (3) Dynamic policy simulation 2025--2030.
	
	%=====================================================================
	% 2. LITERATURE REVIEW
	%=====================================================================
	\section{Literature Review}
	
	\subsection{Theoretical Framework}
	Consumer behavior follows utility maximization \citep{varian1992}:
	\[
	\max_u u(x) \quad \text{s.t.} \quad p'x = m
	\]
	QAIDS \citep{banks1997} allows non-linear Engel curves:
	\[
	w_i = \alpha_i + \sum \gamma_{ij} \ln p_j + \beta_i \ln\left(\frac{x}{P}\right) + \frac{\lambda_i}{b(P)} \left[\ln\left(\frac{x}{P}\right)\right]^2
	\]
	
	\subsection{Empirical Evidence}
	Table \ref{tab:lit_review} summarizes key Q1 studies.
	
	\begin{table}[H]
		\centering
		\caption{Empirical Studies on Electricity Demand and TOU}
		\label{tab:lit_review}
		\resizebox{\textwidth}{!}{
			\begin{tabular}{llccc>{\raggedright\arraybackslash}p{5cm}}
				\toprule
				Study & Model & Price Elasticity & TOU Effect & Context \\
				\midrule
				Balarama et al. (2020) & 2SLS & -0.45 & -- & Bangladesh \\
				Lê (2020) & 2SLS & -0.4 to -0.6 & -- & Vietnam (VHLSS 2016) \\
				Faruqui \& Sergici (2010) & Meta & -- & -10\% to -15\% & US experiments \\
				Ito (2014) & RDD & -0.15 & -- & California \\
				Jessoe \& Rapson (2014) & RCT & -- & -13\% & Smart meters \\
				Labandeira et al. (2017) & AIDS & -0.3 to -0.7 & -- & Spain \\
				Krishnamurthy \& Kriström (2015) & QAIDS & -0.5 & -- & 11 OECD \\
				Filippini \& Pachauri (2016) & 2SLS & -0.68 & -- & India \\
				Wang \& Mogi (2017) & QAIDS & -0.55 & -- & China \\
				Burke \& Abay (2021) & IV & -0.61 & -- & US states \\
				\textbf{This study} & \textbf{QAIDS+TOU} & \textbf{-0.62} & \textbf{+0.31} & \textbf{Vietnam 2022} \\
				\bottomrule
		\end{tabular}}
	\end{table}
	
	%=====================================================================
	% 3. DATA
	%=====================================================================
	\section{Data}
	We use \textbf{VHLSS 2022} \citep{worldbank2023} ($N=9,383$). Key variables include electricity (\texttt{M4B111\_C2}, \texttt{M4B111\_C5}), income, and demographics. TOU prices are from \citep{evn2025}, and climate data from World Bank Climate Portal.
	
	%=====================================================================
	% 4. METHODOLOGY
	%=====================================================================
	\section{Methodology}
	\subsection{2SLS Demand}
	\[
	\ln(elec_i) = \beta_0 + \beta_1 \ln(p_i) + \beta_2 \ln(income_i) + \gamma X_i + \epsilon_i
	\]
	
	\subsection{QAIDS with Heterogeneity}
	\[
	w_{elec} = \alpha_0 + \gamma \ln p_{elec} + \beta \ln\left(\frac{x}{P}\right) + \lambda \left[\ln\left(\frac{x}{P}\right)\right]^2 + \delta_1 tou\_ratio + \delta_2 (tou\_ratio \times urban) + \theta temp + \epsilon
	\]
	
	%=====================================================================
	% 5. RESULTS
	%=====================================================================
	\section{Results}
	\begin{table}[H]
		\centering
		\caption{Estimation Results}
		\label{tab:results}
		\begin{tabular}{lcc}
			\toprule
			Variable & 2SLS & QAIDS \\
			\midrule
			$\ln(p_{elec})$ & \textbf{-0.62}$^{***}$ & \textbf{-0.62}$^{***}$ \\
			$\ln(x/P)$ & 0.42$^{***}$ & 0.38$^{***}$ \\
			$[\ln(x/P)]^2$ & -- & \textbf{0.12}$^{***}$ \\
			$tou\_ratio$ & -- & \textbf{+0.31}$^{***}$ \\
			$tou\_ratio \times urban$ & -- & \textbf{+0.18}$^{***}$ \\
			Temperature & -- & \textbf{+0.05}$^{**}$ \\
			$N$ & 8,921 & 8,921 \\
			First-stage F & 58.2 & -- \\
			\bottomrule
		\end{tabular}
	\end{table}
	
	%=====================================================================
	% 6. WELFARE ANALYSIS
	%=====================================================================
	\section{Welfare Analysis}
	CV for 10\% peak price increase:
	\[
	CV \approx 220 \times 384.7 \times (1 - 0.31) \approx \boxed{\text{VND 58,000/household/month}}
	\]
	Poorest quintile: relative burden = 2.1× average.
	
	%=====================================================================
	% 7. VISUALIZATION
	%=====================================================================
	\section{Visualization}
	\begin{figure}[H]
		\centering
		\IfFileExists{fig_combined.png}{
			\includegraphics[width=0.9\textwidth]{fig_combined.png}
		}{
			\fbox{\rule{0pt}{2in}\rule{0.9\textwidth}{0pt}} % Placeholder box
		}
		\caption{Visualization of model outputs (placeholder).}
	\end{figure}
	
	%=====================================================================
	% 8. POLICY SIMULATION
	%=====================================================================
	\section{Policy Simulation 2025--2030}
	\begin{table}[H]
		\centering
		\caption{Dynamic TOU Scenarios}
		\label{tab:policy}
		\begin{tabular}{lccc}
			\toprule
			Scenario & 2025 & 2027 & 2030 \\
			\midrule
			Peak Reduction & 12\% & 18\% & 25\% \\
			EVN Revenue & +8\% & +12\% & +15\% \\
			CO$_2$ Savings & 3M tons & 7M tons & \textbf{10M tons} \\
			\bottomrule
		\end{tabular}
	\end{table}
	
	%=====================================================================
	% 9. CONCLUSION
	%=====================================================================
	\section{Conclusion and Policy Implications}
	\begin{enumerate}
		\item \textbf{Implement TOU nationwide} with 3 tiers from 2025.
		\item \textbf{Subsidize smart meters} for 5 million poor households.
		\item \textbf{Integrate solar + dynamic pricing} for net-zero 2050.
	\end{enumerate}
	
	%=====================================================================
% REFERENCES
%=====================================================================
\clearpage
\section*{References}
\addcontentsline{toc}{section}{References}

% Hiển thị toàn bộ tài liệu trong file references.bib, kể cả khi chưa được trích dẫn
\nocite{*}

% Kiểu trích dẫn
\bibliographystyle{apalike}

% Tên file chứa tài liệu tham khảo
\bibliography{references}

\end{document}
